\chapter{Mécanismes réactionnels~: les flèches courbes}\label{ch:g12:curly-arrows}

Une \omterm{def:g8:balanced-equations:equation}{équation ajustée} ressemble à la première et à la dernière image d'un
film~: elle montre les \omterm{def:g7:chemical-reactions:reactant}{réactifs} avant et les produits après, et rien
entre les deux. Or les \omterm{def:g7:atoms-and-molecules:molecule}{molécules} ne se contentent pas d'échanger des
\omterm{def:g7:atoms-and-molecules:atom}{atomes}. Les liaisons se rompent et se forment une à une, à mesure que
des doublets d'\omterm{def:g9:inside-the-atom:nucleus}{électrons} se déplacent d'un endroit à un autre, en
passant par des espèces éphémères qui n'apparaissent jamais dans
l'équation. Les chimistes dessinent ces mouvements avec des \omterm{def:g12:curly-arrows:curly-arrow}{flèches
courbes}. Avec quelques règles, ces flèches expliquent pourquoi une
réaction se produit là où elle se produit, et permettent de prévoir des
réactions jamais observées.

\begin{recall}
L'\omterm{def:g11:polarity-and-cohesion:electronegativity}{électronégativité}, les \omterm{def:g11:polarity-and-cohesion:polar-bond}{liaisons polarisées} et les \omterm{def:g11:polarity-and-cohesion:polar-bond}{charges partielles}
(\cref{ch:g11:polarity-and-cohesion}). Les \omterm{def:g10:lewis-and-shape:lewis-structure}{schémas de Lewis}~: \omterm{def:g10:lewis-and-shape:lone-pair}{doublets
liants} et non liants (\cref{ch:g10:lewis-and-shape}). Les \omterm{def:g11:functional-groups:functional-group}{groupes
caractéristiques} et les familles (\cref{ch:g11:functional-groups}).
\end{recall}

\section{Sites riches et sites pauvres en électrons}

\begin{definition}[Sites donneurs et accepteurs d'électrons]\label{def:g12:curly-arrows:site}
Un \emph{site donneur d'électrons}\index{site donneur d'electrons@site donneur d'électrons} d'une
\omterm{def:g7:atoms-and-molecules:molecule}{molécule} ou d'un \omterm{def:g9:ions:ion}{ion} est un endroit riche en \omterm{def:g9:inside-the-atom:nucleus}{électrons}, capable de
céder un doublet pour former une nouvelle liaison~: un \omterm{def:g10:lewis-and-shape:lone-pair}{doublet non liant},
une \omterm{def:g10:lewis-and-shape:multiple-bond}{double liaison}, un \omterm{def:g7:atoms-and-molecules:atom}{atome} portant une charge négative ou une \omterm{def:g11:polarity-and-cohesion:polar-bond}{charge
partielle} $\delta^-$. Un \emph{site accepteur
d'électrons}\index{site accepteur d'electrons@site accepteur d'électrons} est un endroit pauvre en
\omterm{def:g9:inside-the-atom:nucleus}{électrons}, capable de recevoir un doublet~: un \omterm{def:g7:atoms-and-molecules:atom}{atome} portant une charge
positive ou une \omterm{def:g11:polarity-and-cohesion:polar-bond}{charge partielle}
$\delta^+$.
\end{definition}

\begin{example}[Les sites de deux molécules]\label{ex:g12:curly-arrows:sites}
Dans le bromoéthane, \ce{CH3-CH2-Br}, le brome est plus électronégatif
que le carbone~: le carbone qui lui est lié porte $\delta^+$ (site
accepteur), le brome porte $\delta^-$ et trois \omterm{def:g10:lewis-and-shape:lone-pair}{doublets non liants}. Dans
la propanone, l'oxygène du groupe \ce{C=O} porte $\delta^-$ et deux
\omterm{def:g10:lewis-and-shape:lone-pair}{doublets non liants} (sites donneurs), son carbone $\delta^+$ (site
accepteur). L'\omterm{def:g9:ions:ion}{ion} hydroxyde \ce{OH-}, avec sa charge négative et ses
trois \omterm{def:g10:lewis-and-shape:lone-pair}{doublets non liants}, est un donneur puissant.
\end{example}

\begin{omfigure}
\setchemfig{atom sep=2.2em}
\begin{tabular}{cc}
\chemfig{H_3C-C(-[2]H)(-[6]H)-\charge{90:2pt=\:,0=\:,270:2pt=\:}{Br}} &
\chemfig{H_3C-C(-[6]CH_3)=[2]\charge{0=\:,180=\:}{O}} \\[2pt]
\small $\delta^+$ sur le carbone lié à \ce{Br}, $\delta^-$ sur \ce{Br} &
\small $\delta^+$ sur le carbone de \ce{C=O}, $\delta^-$ sur \ce{O} \\
\small bromoéthane & \small propanone
\end{tabular}

{\small Des \omterm{def:g7:atoms-and-molecules:atom}{atomes} de carbone pauvres en \omterm{def:g9:inside-the-atom:nucleus}{électrons} ($\delta^+$, sites
accepteurs) voisins d'\omterm{def:g7:atoms-and-molecules:atom}{atomes} électronégatifs riches en \omterm{def:g9:inside-the-atom:nucleus}{électrons}
($\delta^-$, porteurs de \omterm{def:g10:lewis-and-shape:lone-pair}{doublets non liants}~: sites donneurs).}
\end{omfigure}

\section{Les flèches courbes}

\begin{definition}[Flèche courbe]\label{def:g12:curly-arrows:curly-arrow}
Une \emph{flèche courbe}\index{fleche courbe@flèche courbe} représente le mouvement
d'un doublet d'\omterm{def:g9:inside-the-atom:nucleus}{électrons} au cours d'une étape d'une réaction. Elle part
du doublet qui se déplace, \omterm{def:g10:lewis-and-shape:lone-pair}{doublet non liant} ou liaison, et aboutit là
où va le doublet~: sur un \omterm{def:g7:atoms-and-molecules:atom}{atome}, avec lequel il forme une nouvelle
liaison, ou sur un \omterm{def:g7:atoms-and-molecules:atom}{atome} qui récupère le doublet d'une liaison qui se
rompt.
\end{definition}

\begin{method}[Dessiner des flèches courbes]\label{met:g12:curly-arrows:drawing}
\begin{enumerate}
  \item Repérer le site donneur et le site accepteur.
  \item Dessiner une flèche du doublet donneur (\omterm{def:g10:lewis-and-shape:lone-pair}{doublet non liant} ou
        liaison) vers l'\omterm{def:g7:atoms-and-molecules:atom}{atome} accepteur~: une nouvelle liaison s'y forme.
  \item Si l'\omterm{def:g7:atoms-and-molecules:atom}{atome} accepteur se retrouve alors avec trop d'\omterm{def:g9:inside-the-atom:nucleus}{électrons}
        (plus d'un octet, ou plus de deux pour l'hydrogène), dessiner une
        seconde flèche qui rompt l'une de ses liaisons, le doublet allant
        vers l'\omterm{def:g7:atoms-and-molecules:atom}{atome} le plus électronégatif.
  \item Vérifier les charges après l'étape~: l'\omterm{def:g7:atoms-and-molecules:atom}{atome} qui a cédé un
        \omterm{def:g10:lewis-and-shape:lone-pair}{doublet non liant} gagne une charge positive, l'\omterm{def:g7:atoms-and-molecules:atom}{atome} qui a reçu
        un doublet sous forme de \omterm{def:g10:lewis-and-shape:lone-pair}{doublet non liant} gagne une charge
        négative~; la charge totale ne change pas.
\end{enumerate}
\end{method}

\begin{proposition}[La charge se conserve à chaque étape]\label{prop:g12:curly-arrows:charge}
La charge électrique totale des espèces est la même avant et après
chaque étape d'un mécanisme.
\end{proposition}

\begin{proof}
Une \omterm{def:g12:curly-arrows:curly-arrow}{flèche courbe} ne fait que déplacer des \omterm{def:g9:inside-the-atom:nucleus}{électrons} déjà présents~;
aucun \omterm{def:g9:inside-the-atom:nucleus}{électron} n'est créé ni détruit, et les noyaux ne changent pas.
\end{proof}

\begin{example}[Une première flèche]\label{ex:g12:curly-arrows:first}
Un \omterm{def:g9:ions:ion}{ion} hydrogène \ce{H+} rencontre un \omterm{def:g9:ions:ion}{ion} hydroxyde \ce{OH-}. Un \omterm{def:g10:lewis-and-shape:lone-pair}{doublet
non liant} de l'oxygène forme une liaison avec l'\omterm{def:g9:ions:ion}{ion} hydrogène, qui n'a
aucun \omterm{def:g9:inside-the-atom:nucleus}{électron}~:
\ce{H+ + OH- -> H2O}. Une seule flèche, du \omterm{def:g10:lewis-and-shape:lone-pair}{doublet non liant} de
l'oxygène vers \ce{H+}~; les charges $+1$ et $-1$ deviennent 0 dans la
\omterm{def:g7:atoms-and-molecules:molecule}{molécule} d'eau.
\end{example}

\begin{omfigure}
\setchemfig{atom sep=2.1em}
\chemfig{@{hp}H^{+}}\qquad\raisebox{0.2ex}{$+$}\qquad
\chemfig{H-@{ox}\charge{90=\:,0=\:,270=\:}{O}^{-}}
\qquad\raisebox{0.2ex}{$\longrightarrow$}\qquad
\chemfig{H-\charge{90=\:,270=\:}{O}-H}
\chemmove{
  \draw[->, thick, omProp, shorten <=4pt, shorten >=3pt]
    ($(ox)+(0,0.25)$) .. controls +(110:9mm) and +(60:9mm) .. (hp);}

{\small L'étape la plus simple~: un \omterm{def:g10:lewis-and-shape:lone-pair}{doublet non liant} de l'\omterm{def:g9:ions:ion}{ion}
hydroxyde se déplace pour former la nouvelle liaison \ce{O-H} avec l'\omterm{def:g9:ions:ion}{ion}
hydrogène.}
\end{omfigure}

\section{Actes élémentaires et intermédiaires réactionnels}

\begin{definition}[Mécanisme, acte élémentaire, intermédiaire]\label{def:g12:curly-arrows:mechanism}
Le \emph{mécanisme réactionnel}\index{mecanisme reactionnel@mécanisme réactionnel} d'une
réaction est la suite des étapes par lesquelles les \omterm{def:g7:chemical-reactions:reactant}{réactifs} deviennent
les produits. Chaque \emph{acte élémentaire}\index{acte elementaire@acte élémentaire} est
un événement unique au cours duquel quelques doublets d'\omterm{def:g9:inside-the-atom:nucleus}{électrons} se
déplacent à la fois. Une espèce formée lors d'une étape et consommée
lors d'une étape ultérieure est un \emph{intermédiaire
réactionnel}\index{intermediaire reactionnel@intermédiaire réactionnel}~:
elle n'apparaît pas dans l'équation globale.
\end{definition}


\begin{example}[Un carbocation intermédiaire]\label{ex:g12:curly-arrows:carbocation}
Quand le bromure d'hydrogène s'additionne sur l'éthène, la réaction se
fait en deux étapes. Dans la première, la \omterm{def:g10:lewis-and-shape:multiple-bond}{double liaison} de l'éthène
cède un doublet à l'hydrogène de
\ce{HBr}, et le doublet de \ce{H-Br} va vers le brome, qui part sous
forme de
\ce{Br-}~; le carbone qui a perdu sa part de la \omterm{def:g10:lewis-and-shape:multiple-bond}{double liaison} ne garde
que six \omterm{def:g9:inside-the-atom:nucleus}{électrons} et porte une charge positive~: c'est un carbocation,
\ce{CH3-CH2+}. Dans la seconde étape, un \omterm{def:g10:lewis-and-shape:lone-pair}{doublet non liant} de \ce{Br-}
forme une liaison avec ce carbone. Le carbocation est un intermédiaire.
\end{example}

\section{Trois catégories de réactions}

\begin{definition}[Substitution, addition, élimination]\label{def:g12:curly-arrows:categories}
Lors d'une \emph{substitution}\index{substitution}, un \omterm{def:g7:atoms-and-molecules:atom}{atome} ou un
groupe d'\omterm{def:g7:atoms-and-molecules:atom}{atomes} d'une \omterm{def:g7:atoms-and-molecules:molecule}{molécule} est remplacé par un autre. Lors d'une
\emph{addition}\index{addition}, deux \omterm{def:g7:atoms-and-molecules:molecule}{molécules} s'unissent en une seule,
une liaison multiple devenant simple. Lors d'une
\emph{élimination}\index{elimination@élimination}, une petite \omterm{def:g7:atoms-and-molecules:molecule}{molécule} (souvent
l'eau) est retirée d'une \omterm{def:g7:atoms-and-molecules:molecule}{molécule}, une liaison simple devenant multiple.
\end{definition}

\begin{omfigure}
\setchemfig{atom sep=2em}
\small
\textbf{\omterm{def:g12:curly-arrows:categories}{Substitution}} (une étape)~: \ce{CH3-CH2-Br + OH- -> CH3-CH2-OH + Br-}

\medskip
\chemfig{H-@{o1}\charge{90=\:,0=\:,270=\:}{O}^{-}}\qquad
\chemfig{H_3C-@{c1}C(-[2]H)(-[6]H)-[@{b1}]@{br1}\charge{90=\:,0=\:,270=\:}{Br}}
\qquad$\longrightarrow$\qquad
\chemfig{H_3C-C(-[2]H)(-[6]H)-O-H}\quad$+$\quad\ce{Br-}
\chemmove{
  \draw[->, thick, omProp, shorten <=4pt, shorten >=3pt]
    ($(o1)+(0.15,0.2)$) .. controls +(60:9mm) and +(120:9mm) .. (c1);
  \draw[->, thick, omDef, shorten <=2pt, shorten >=4pt]
    (b1) .. controls +(-90:7mm) and +(-90:7mm) .. (br1);}

\vspace{6mm}
\textbf{\omterm{def:g12:curly-arrows:categories}{Addition}} (deux étapes)~: \ce{CH2=CH2 + HBr -> CH3-CH2-Br}

\vspace{9mm}
\chemfig{H_2C=[@{pi}]CH_2}\quad$+$\quad
\chemfig{@{h2}H-[@{b2}]@{br2}Br}
\qquad$\longrightarrow$\qquad
\chemfig{H_3C-@{cp}\charge{90:2pt=$\scriptstyle +$}{C}H_2}\quad$+$\quad
\chemfig{@{brm}\charge{90=\:,180=\:,270=\:,0=\:}{Br}^{-}}
\qquad$\longrightarrow$\qquad \ce{CH3-CH2-Br}
\chemmove{
  \draw[->, thick, omProp, shorten <=2pt, shorten >=3pt]
    (pi) .. controls +(90:10mm) and +(110:9mm) .. (h2);
  \draw[->, thick, omDef, shorten <=2pt, shorten >=4pt]
    (b2) .. controls +(-90:7mm) and +(-90:7mm) .. (br2);
  \draw[->, thick, omProp, shorten <=4pt, shorten >=3pt]
    ($(brm)+(0,-0.25)$) .. controls +(-120:8mm) and +(-60:8mm) .. (cp);}

\vspace{6mm}
\textbf{\omterm{def:g12:curly-arrows:categories}{Élimination}} (déshydratation de l'éthanol en milieu acide, trois étapes)~:
\ce{CH3-CH2-OH -> CH2=CH2 + H2O}

\medskip
\begin{tabular}{@{}l@{}}
1.\ l'oxygène capte un \omterm{def:g9:ions:ion}{ion} hydrogène~: \ce{CH3-CH2-OH + H+ -> CH3-CH2-OH2+}~;\\
2.\ le doublet \ce{C-O} part avec l'eau~: \ce{CH3-CH2-OH2+ -> CH3-CH2+ + H2O}~;\\
3.\ un doublet \ce{C-H} voisin devient la \omterm{def:g10:lewis-and-shape:multiple-bond}{double liaison}, et
    \ce{H+} part~:\\
\phantom{3.\ }\ce{CH3-CH2+ -> CH2=CH2 + H+}.
\end{tabular}

\medskip
{\small Trois mécanismes. Flèches orange~: un doublet donneur forme une
nouvelle liaison. Flèches bleues~: une liaison se rompt, son doublet
allant vers l'\omterm{def:g7:atoms-and-molecules:atom}{atome} le plus électronégatif. Dans l'\omterm{def:g12:curly-arrows:categories}{élimination}, l'\omterm{def:g9:ions:ion}{ion}
hydrogène est restitué à la fin.}
\end{omfigure}


\begin{remark}[Modifier le squelette ou le groupe]\label{rem:g12:curly-arrows:skeleton}
Certaines réactions ne modifient que le \omterm{def:g11:functional-groups:functional-group}{groupe caractéristique} et
conservent le squelette carboné~: la \omterm{def:g12:curly-arrows:categories}{substitution} du bromoéthane
transforme un \omterm{def:g11:functional-groups:halogenoalkane}{halogénoalcane} en \omterm{def:g11:functional-groups:alcohol}{alcool}, avec deux carbones avant et deux
après. D'autres construisent ou coupent le squelette~: former une
nouvelle liaison carbone--carbone allonge une chaîne, l'étape clé pour
fabriquer une grosse \omterm{def:g7:atoms-and-molecules:molecule}{molécule} à partir de petites. Planifier une
\omterm{def:g10:synthesis-yield:synthesis}{synthèse}, c'est choisir des réactions des deux sortes, dans le bon
ordre.
\end{remark}

\section{Exercices}

\begin{exercise}[$\star$]\label{exo:g12:curly-arrows:1}
Pour chaque espèce, nommer un site donneur et, s'il y en a un, un site
accepteur~: \ce{OH-}~; \ce{H+}~; \ce{CH2=CH2}~; \ce{CH3-Cl}~; \ce{H2O}.
\end{exercise}

\begin{exercise}[$\star$]\label{exo:g12:curly-arrows:2}
\omterm{def:g12:curly-arrows:categories}{Substitution}, \omterm{def:g12:curly-arrows:categories}{addition} ou \omterm{def:g12:curly-arrows:categories}{élimination}~?
(a) \ce{CH3-CH2-Cl + OH- -> CH3-CH2-OH + Cl-};
(b) \ce{CH2=CH2 + H2O -> CH3-CH2-OH};
(c) \ce{CH3-CH2-OH -> CH2=CH2 + H2O};
(d) \ce{CH2=CH2 + Br2 -> CH2Br-CH2Br}.
\end{exercise}

\begin{exercise}[$\star$]\label{exo:g12:curly-arrows:3}
D'où part une \omterm{def:g12:curly-arrows:curly-arrow}{flèche courbe}, et où aboutit-elle~? Que signifie une
\omterm{def:g12:curly-arrows:curly-arrow}{flèche courbe} allant d'un \omterm{def:g10:lewis-and-shape:lone-pair}{doublet non liant} d'un \omterm{def:g7:atoms-and-molecules:atom}{atome} d'oxygène vers un
\omterm{def:g7:atoms-and-molecules:atom}{atome} de carbone~?
\end{exercise}

\begin{exercise}[$\star$]\label{exo:g12:curly-arrows:4}
Dans la réaction \ce{H+ + NH3 -> NH4+}, quelle espèce cède le doublet~?
Dessiner la flèche. Quelle est la charge du produit~?
\end{exercise}

\begin{exercise}[$\star$]\label{exo:g12:curly-arrows:5}
Qu'est-ce qu'un \omterm{def:g12:curly-arrows:mechanism}{intermédiaire réactionnel}~? Nommer l'intermédiaire de
l'\omterm{def:g12:curly-arrows:categories}{addition} de
\ce{HBr} sur l'éthène.
\end{exercise}

\begin{exercise}[$\star\star$]\label{exo:g12:curly-arrows:6}
Dessiner les \omterm{def:g12:curly-arrows:curly-arrow}{flèches courbes} de la \omterm{def:g12:curly-arrows:categories}{substitution}
\ce{CH3-Br + OH- -> CH3-OH + Br-} (une étape), et vérifier les charges.
\end{exercise}

\begin{exercise}[$\star\star$]\label{exo:g12:curly-arrows:7}
Lors de la première étape de l'\omterm{def:g12:curly-arrows:categories}{addition} de \ce{HBr} sur l'éthène,
pourquoi le doublet de
\ce{H-Br} va-t-il vers le brome plutôt que vers l'hydrogène~? Quelle
charge porte chaque produit de cette étape~?
\end{exercise}

\begin{exercise}[$\star\star$]\label{exo:g12:curly-arrows:8}
Pour la déshydratation de l'éthanol, écrire les trois étapes et dire,
pour chacune, quelle espèce est le donneur et laquelle l'accepteur.
Pourquoi l'\omterm{def:g9:ions:ion}{ion} hydrogène n'apparaît-il pas dans l'équation globale~?
\end{exercise}

\begin{exercise}[$\star\star$]\label{exo:g12:curly-arrows:9}
L'étape \ce{CH3-CH2+ + H2O -> CH3-CH2-OH2+} suit
l'\omterm{def:g12:curly-arrows:categories}{addition} d'un \omterm{def:g9:ions:ion}{ion} hydrogène sur l'éthène dans l'eau. Dessiner la
flèche. Qu'obtient-on après une nouvelle perte de \ce{H+}~? Écrire
l'équation globale et donner sa catégorie.
\end{exercise}

\begin{exercise}[$\star\star$]\label{exo:g12:curly-arrows:10}
Lire le mécanisme de la \omterm{def:g12:curly-arrows:categories}{substitution} du bromoéthane~: combien de
doublets se déplacent~? Quelle liaison se forme, laquelle se rompt~?
Y a-t-il un intermédiaire~?
\end{exercise}

\begin{exercise}[$\star\star$]\label{exo:g12:curly-arrows:11}
% ledger: en:C, en:Cl, en:Br
Le chlorométhane réagit avec l'\omterm{def:g9:ions:ion}{ion} hydroxyde plus lentement que le
bromométhane. À l'aide des \omterm{def:g11:polarity-and-cohesion:electronegativity}{électronégativités} (C 2.55, Cl 3.16, Br 2.96),
expliquer quel carbone est le plus $\delta^+$. Cela explique-t-il les
vitesses~? (Penser aussi à la facilité avec laquelle l'\omterm{def:g9:ions:ion}{ion} partant
emporte le doublet.)
\end{exercise}

\begin{exercise}[$\star\star\star$]\label{exo:g12:curly-arrows:12}
Le chlorure d'hydrogène s'additionne sur le propène, \ce{CH2=CH-CH3}.
Lors de la première étape, l'\omterm{def:g9:ions:ion}{ion} hydrogène peut se lier à l'un ou
l'autre carbone de la \omterm{def:g10:lewis-and-shape:multiple-bond}{double liaison}, ce qui donne deux carbocations
possibles. Les dessiner tous les deux, ainsi que les deux produits
possibles, et les nommer. (Lequel se forme majoritairement est étudié
dans le volume de première année d'université.)
\end{exercise}

\begin{exercise}[$\star\star\star$]\label{exo:g12:curly-arrows:13}
Un élève dessine la première étape de la \omterm{def:g12:curly-arrows:categories}{substitution} du bromoéthane
avec une flèche allant de l'\omterm{def:g7:atoms-and-molecules:atom}{atome} de brome vers l'oxygène de \ce{OH-}.
Expliquer les deux erreurs.
\end{exercise}

\begin{exercise}[$\star\star\star$]\label{exo:g12:curly-arrows:14}
Lors de la formation d'un \omterm{def:g11:functional-groups:carboxylic-acid}{ester} à partir d'un \omterm{def:g11:functional-groups:carboxylic-acid}{acide carboxylique} et d'un
\omterm{def:g11:functional-groups:alcohol}{alcool}, l'oxygène de l'\omterm{def:g11:functional-groups:alcohol}{alcool} attaque le carbone du groupe \ce{C=O}. À
l'aide des \omterm{def:g11:polarity-and-cohesion:polar-bond}{charges partielles}, expliquer pourquoi ce carbone est un site
accepteur et cet oxygène un site donneur. Quelle petite \omterm{def:g7:atoms-and-molecules:molecule}{molécule} est
perdue au total, et à quelle catégorie appartient la réaction dans son
ensemble~?
\end{exercise}

\begin{exercise}[$\star\star\star$]\label{exo:g12:curly-arrows:15}
L'éthène réagit avec le dibrome \ce{Br2}, une \omterm{def:g11:polarity-and-cohesion:polar-molecule}{molécule apolaire}. Quand
une \omterm{def:g7:atoms-and-molecules:molecule}{molécule} de \ce{Br2} s'approche de la \omterm{def:g10:lewis-and-shape:multiple-bond}{double liaison}, les \omterm{def:g9:inside-the-atom:nucleus}{électrons}
de \ce{Br-Br} sont repoussés loin de l'\omterm{def:g7:atoms-and-molecules:atom}{atome} de brome le plus proche.
Expliquer pourquoi cela crée un site accepteur, et proposer la première
étape du mécanisme.
\end{exercise}

\section{Problème~: fabriquer du bromoéthane}

\begin{problem}[{Problème du week-end --- à partir de dix grammes
d'éthène~: combien de bromoéthane~?}]
\label{pb:g12:curly-arrows:1}
% ledger: aw:C, aw:H, aw:Br, en:H, en:Br, en:C
Le bromoéthane, \omterm{def:g7:chemical-reactions:reactant}{réactif} de départ de nombreuses \omterm{def:g10:synthesis-yield:synthesis}{synthèses}, est préparé
par \omterm{def:g12:curly-arrows:categories}{addition} de bromure d'hydrogène sur l'éthène~:
\ce{CH2=CH2 + HBr -> CH3-CH2-Br}. Un chimiste
utilise \qty{10.0}{g} d'éthène et un excès de bromure d'hydrogène~; le
\omterm{def:g10:synthesis-yield:yield}{rendement} est de \qty{75}{\percent}.

\medskip
\noindent\textbf{Partie I --- Les sites.}
\begin{enumerate}
  \item Dessiner les \omterm{def:g10:lewis-and-shape:lewis-structure}{schémas de Lewis} de l'éthène et du bromure
        d'hydrogène.
  \item Où se trouve le site donneur de l'éthène~? Pourquoi~?
  \item Calculer la différence d'\omterm{def:g11:polarity-and-cohesion:electronegativity}{électronégativité} de la liaison
        \ce{H-Br}
        (H 2.20, Br 2.96). Quel \omterm{def:g7:atoms-and-molecules:atom}{atome} porte $\delta^+$~?
  \item Quel \omterm{def:g7:atoms-and-molecules:atom}{atome} de \ce{HBr} est le site accepteur~?
\end{enumerate}

\medskip
\noindent\textbf{Partie II --- Le mécanisme.}
\begin{enumerate}[resume]
  \item Dessiner la première étape avec ses deux \omterm{def:g12:curly-arrows:curly-arrow}{flèches courbes}.
  \item Donner les deux espèces formées et leurs charges.
  \item Vérifier la conservation de la charge lors de la première étape.
  \item Dessiner la seconde étape avec sa \omterm{def:g12:curly-arrows:curly-arrow}{flèche courbe}.
  \item Combien d'\omterm{def:g9:inside-the-atom:nucleus}{électrons} entourent le carbone positif de
        l'intermédiaire~? Pourquoi est-il si \omterm{def:g7:chemical-reactions:reactant}{réactif}~?
\end{enumerate}

\medskip
\noindent\textbf{Partie III --- Quel type de réaction~?}
\begin{enumerate}[resume]
  \item Quel est l'intermédiaire du mécanisme~?
  \item Pourquoi n'apparaît-il pas dans l'équation globale~?
  \item À quelle catégorie la réaction appartient-elle~?
  \item La réaction modifie-t-elle le squelette carboné, ou seulement le
        \omterm{def:g11:functional-groups:functional-group}{groupe caractéristique}~?
\end{enumerate}

\medskip
\noindent\textbf{Partie IV --- Quelle quantité de produit~?}
\begin{enumerate}[resume]
  \item Calculer les \omterm{def:g10:the-mole:molar-mass}{masses molaires} de l'éthène et du bromoéthane.
  \item Calculer la quantité d'éthène.
  \item Remplir le \omterm{def:g10:reaction-progress-table:extent}{tableau d'avancement}. Quel \omterm{def:g7:chemical-reactions:reactant}{réactif} est limitant~?
  \item Calculer la masse maximale de bromoéthane.
  \item Calculer la masse obtenue avec un \omterm{def:g10:synthesis-yield:yield}{rendement} de
        \qty{75}{\percent}.
  \item Donner la réponse finale~: quelle masse de bromoéthane le
        chimiste obtient-il~?
\end{enumerate}
\end{problem}
